\section{A semistable log model}\label{sec:log-model}

The algebraic comparison requires a model on which the exceptional curves
of the general surface form a horizontal boundary.  The following construction
keeps that boundary while making the special fibre reduced.  Its output is the
input for the support and adjunction arguments in the next section.

\begin{proposition}[The relative log model]\label{prop:log-model}
Let $C$ be a smooth connected complex algebraic curve with a marked point $0$,
and let $\pi:X\to C$ be a flat projective family of integral normal Gorenstein
surfaces.  Suppose that $X$ is normal and Gorenstein, that a section
$\Gamma\subset X$ contains every fibre singularity, and that $\pi$ is smooth
on $X\setminus\Gamma$.  Then, after shrinking $C$ about $0$ and making a
finite base change from a smooth curve, there is a projective birational
morphism
\[
 h:Y\longrightarrow X
\]
with the following properties.  Here $X$, $C$, $\Gamma$ and $0$ denote the
base-changed objects, and $H$ is the sum, with coefficient one, of all
$h$-exceptional prime divisors dominating $C$.
\begin{enumerate}
\item $Y$ is normal and $\mathbb Q$-factorial, and $h$ is an isomorphism
      over $X\setminus\Gamma$.
\item For every $c\in C$, the scheme fibre $Y_c$ is a reduced Cartier divisor,
      and $(Y,H+Y_c)$ is divisorial log terminal.
\item $K_Y+H$ is nef over $X$.  Every vertical $h$-exceptional prime divisor
      lies in $Y_0$ and maps to the point $\Gamma\cap X_0$.
\item For $c\ne0$, the surface $Y_c$ is integral and normal, and the divisor
      $H|_{Y_c}$ is reduced as a Weil divisor.  The fibre $Y_0$ has a unique component dominating
      $X_0$; its remaining components are projective surfaces over
      $\Gamma\cap X_0$.
\end{enumerate}
\end{proposition}

Here nefness over $X$ means nonnegative intersection with every complete curve
contracted by $h$.  A divisorial log terminal (dlt) pair is used in the usual
characteristic-zero sense; in particular its coefficient-one components are
normal and it has simple normal crossings at the generic points of their
intersection strata.  We construct the marked semistable model by
Proposition~\ref{prop:marked-monomial-reduction}, which makes explicit the
dimension-three monomial reduction of
\cite[Theorems~7.17 and~7.19]{KM1998}, and then apply the semistable minimal
model theorem \cite[Theorem~7.9]{KM1998}.  In its terminology, a dlt morphism
to a smooth curve requires the fibre to be included in the boundary at every
parameter \cite[Definition~7.1]{KM1998}.

\begin{proof}
First choose a projective log resolution $g:Z\to X$ which principalizes the
ideal of $\Gamma$ and resolves the union with $X_0$.  We choose it to be an
isomorphism over $X\setminus\Gamma$: that open set is smooth and its special
fibre is already smooth.  Thus $Z$ is smooth, $g^{-1}(\Gamma)$ has divisorial
support, and the union of this support with the support of $Z_0$ has simple
normal crossings.  All exceptional divisors of $g$ lie over $\Gamma$.
By generic smoothness for $Z\to C$ and for the finitely many strata of its
exceptional divisor, we may delete finitely many base points other than $0$
so that $Z\to C$ is smooth away from $0$ and every fibre support together with
the exceptional divisor has simple normal crossings.  After the same
shrinking, every vertical exceptional divisor lies over $0$.

Let $H_Z$ be the reduced horizontal exceptional divisor of $g$, and write
$Z_0=\sum a_jZ_j$.  The preceding shrinking makes
$(Z_0)_{\mathrm{red}}+H_Z$ simple normal crossing, with $H_Z$ relatively
simple normal crossing over $C\setminus\{0\}$.  Choose a positive integer
$b$ divisible by every $a_j$, and a parameter $\tau$ at $0$ regular with no
other zero on a Zariski neighbourhood.  The cover $s^b=\tau$, with its smooth
normalization, is finite, ramified to order $b$ over $0$, and \'etale
elsewhere.  We again write $X$, $C$ and $\Gamma$ for the pulled-back objects.

The pulled-back target $X$ is still normal and Gorenstein.  Indeed the chosen
curve cover is locally $s^b=\tau$, so the target is locally the hypersurface
$A[s]/(s^b-\tau)$ over a Gorenstein ring $A$.  It is therefore Gorenstein and
satisfies Serre's condition $S_2$.  Off the section it is smooth over the new
smooth curve, while the section has codimension two.  Thus it satisfies
$R_1$ and is normal.  Its fibres remain integral, so the flat total space over
the connected covering curve is integral.

Apply Proposition~\ref{prop:marked-monomial-reduction} to $Z$ with marking
$H_Z$ and the chosen cover.  It gives a smooth variety $V$, projective over
the normalized pullback of $Z$, whose special fibre is reduced and simple
normal crossing together with the transformed horizontal marking.
The construction is an isomorphism on the normalized pullback away from $0$,
where the family and marking already have these properties.  Over
$X\setminus\Gamma$, the original morphism was smooth and the marking was
empty, so the same assertion shows that $V\to X$ is an isomorphism there.
Thus all exceptional divisors still lie over $\Gamma$, and every
vertical exceptional divisor lies over $0$.

Let $H_V$ be the reduced sum of the horizontal exceptional prime divisors on
$V$.  Since the modification is an isomorphism away from $0$, $H_V$ is
precisely the transformed reduced marking $H_Z$.  For every $c$, the pair
$(V,H_V+V_c)$ is simple normal crossing.
Consequently $(V,H_V)\to C$ is a dlt morphism.  The relative minimal model
program of \cite[Theorem~7.9]{KM1998} applies to the threefold pair $(V,H_V)$
over the flat target $X\to C$.  All its contractions and flips are projective
over $X$, it terminates, and the resulting morphism $(Y,H)\to C$ is again
dlt.  Its alternative of a Mori fibre space cannot occur: such a space would
factor the dominant morphism from the threefold to $X$ through a variety of
dimension less than three.  We therefore obtain $K_Y+H$ nef over $X$.

The program consists of divisorial contractions and flips, so it extracts
no divisors.  Every surviving horizontal exceptional prime consequently
retains coefficient one in $H$, and every vertical exceptional prime still
lies over $0$.  Each step is an isomorphism over $X\setminus\Gamma$, since
there were no relative curves to contract there.  The usual preservation of
$\mathbb Q$-factoriality in this program gives the asserted
$\mathbb Q$-factorial model.

For completeness, the scheme-theoretic fibre is precisely the boundary fibre
used above.  The pullback of a local parameter at $c$ is a nonzerodivisor on
the integral normal variety $Y$, so $Y_c$ is Cartier.  Its order at every prime
component is one: the corresponding divisorial valuation occurred on $V$,
where the fibre was reduced, and the program did not extract divisors.
A principal divisor with these orders on a normal variety defines a reduced
subscheme.  Locally, if $u^n\in(t)$, the height-one valuations show that
$u/t$ is regular at every height-one prime; normality then gives $u/t\in
\mathcal O_Y$.  Thus $(t)$ is radical.

For $c\ne0$ there are no vertical exceptional divisors, and $X_c$ is integral.
Its strict transform is therefore the only component of $Y_c$.  This component
is normal by the dlt property.  At the generic point of each curve in
$H\cap Y_c$, the pair is simple normal crossing; hence the Cartier fibre cuts
$H$ with multiplicity one, giving the reduced Weil divisor $H|_{Y_c}$.  Over $0$
there is again just one strict transform of $X_0$; all other components are
exceptional.  They map into $\Gamma\cap X_0$ and are projective, because $h$
is projective.  This proves all the assertions.
\end{proof}

The proposition identifies $H$ as the full surviving horizontal divisorial
boundary.  It does not yet identify the whole set $h^{-1}(\Gamma)$ with the
boundary, nor identify the different in adjunction on each fibre component.
Those are separate consequences of nefness, non-log-canonicity and the support
argument proved next.
